\section{Basic Description}

We shall describe a standard version of type theory called (intensional) Martin-L\"of Type Theory (MLTT~\cite{martin1975intuitionistic}). 
Set theoretic foundations usually consist of a number of \textit{axioms}, written in the language of 
first-order logic, which is separately assumed beforehand. The logic's deductive system provides rules for making  
inferences, while the \textit{axioms} define and describe sets and
their semantics. This forms a basis to build mathematics, using sets to describe mathematical
objects, while manipulating statements via first-order logic.
MLTT limits the dependence on axioms and constructs more machinery within the deductive system itself. The
deductive system consists of the primitive notion of `types', the building blocks
for mathematical objects as well as the logic, while specifying the rule that govern them. 

MLTT consists of \textit{rules} for deriving \textit{judgements}, which are either of the form
$ a : A$ meaning `$a$ is a term/instance of type A', or $a \deq b : A$ meaning `$a$ is judgementally equal to $b$',
while both being instances of $A$.
Judgemental equality is a strict notion of equality which classifies objects as definitionally
the same within the intrinsic language of the formal system.
A separate notion `propositional equality', denoted with $=$, is a weaker notion of equality, corresponding
to the equality of the mathematical objects to be described. This is only a feature of \emph{intensional} MLTT.

The language of MLTT is an extension of dependently-typed $\lambda$-calculus, with variables
$x_1, x_2, \ldots$, primitive constants $c, c', c'', \ldots$, and defined constants $f, f', f''$.
The terms $t, t', t'', \ldots$ are given by the grammar 
\begin{align*}
  t ::= x \mid \lambda x.t \mid t(t') \mid  c \mid f.
\end{align*}

Further, MLTT provides rules to form types from a preexisting context, and rules to manipulate defined or preexisting types 
to encode their semantics.

The logic of MLTT comes encoding propositions as types, wherein having an term of a `proposition as a type'
is interpreted as a witness or evidence of the proposition.~\cite{awodey2004propositions} MLTT being an intuitionistic formal system, has a constructive
logic, which requires any existence proof to be certified by a witness. ~\cite{martin1975intuitionistic}

\subsection{Universes}

Type Theory `postulates' a chain of cumulative universes à la Russell,
  $\U_0, \U_1, \U_2, \ldots$ as primitive constants
with
\begin{itemize}
  \item $\U_i : \U_{i+1}$ for each $i$ 
  \item If $A : \U_i$, then $A : \U_{i+1}$ 
\end{itemize}
Any type in the theory is an instance of some universe in this chain. Each universe is also a type,
and is an instance of universes above it in this heirarchy. Unless necessary, we shall often omit the index from
the universe and write $A : \U$.

\textit{\textbf{Remark.} This chain, while indexed by natural numbers, is not related to the Natural 
Numbers that will be defined in the theory. The indices come the metatheory.
The arithmetic object of Natural Numbers $(\N)$ is richer and needs more
information to be properly defined.}
